% GENERATED FILE. Edit canonical lessons, then run npm run generate:books.
% canonical-source-hash: fnv1a64:316fba419480f840
% canonical-lessons: HI-C330-ikattha, HI-C330-jhaki, HI-C330-jhulna, HI-C330-sobhayatra, HI-C330-bhakt

\chapter{At the Fair}
\label{ch:hi-at-the-fair}
\hlchaptermodality{330}

\begin{chapteropening}
\textbf{By the end of this chapter:} \emph{I can describe the crowds and the sights at a fair.}

Complete the last lesson of chapter 330: I can describe the crowds and the sights at a fair.
\end{chapteropening}

\section[\texorpdfstring{\hi{इकट्ठा} (ikaṭṭhā)}{इकट्ठा (ikaṭṭhā)}]{\hi{इकट्ठा} (ikaṭṭhā) — gathered together}

\label{lesson:HI-C330-ikattha}



\begin{quote}
Before the new one: say the Hindi for to sway (with joy or music), then the Hindi for merriment.
\end{quote}

\subsection*{You'll want to know: \hi{इकट्ठा}}
\textbf{\hi{इकट्ठा}} — \emph{ikaṭṭhā} — \textquotedblleft{}gathered together\textquotedblright{}.

\begin{quote}
\hi{मेले} \hi{में} \hi{पूरा} \hi{गाँव} \hi{इकट्ठा} \hi{हुआ।}
\end{quote}

\begin{quote}
\emph{mele meṁ pūrā gā̃v ikaṭṭhā huā.}
\end{quote}

\begin{quote}
The whole village gathered at the fair.
\end{quote}

\begin{cousinweb}
\textbf{\hi{इकट्ठा}} (\emph{ikaṭṭhā}) goes back to Sanskrit \emph{ekastha}, \textquotedblleft{}standing together\textquotedblright{} (CDIAL 2480): all in one spot. \textbf{\hi{इकट्ठा} \hi{होना}} is to gather; \textbf{\hi{इकट्ठा} \hi{करना}}, to collect.
\end{cousinweb}

\begin{grammarlens}[title={What you've built}]
One more word for telling what happened.
\end{grammarlens}

\subsection*{Your turn}
\begin{itemize}
\raggedright
  \item \emph{Say it:} \emph{ikaṭṭhā}
  \item \emph{Say it:} \emph{ikaṭṭhā}, once more
  \item \emph{Say it:} \emph{mele meṁ pūrā gā̃v ikaṭṭhā huā.} — The whole village gathered at the fair.
  \item \emph{Recall:} say the Hindi for to sway (with joy or music), then the Hindi for merriment, then say \emph{ikaṭṭhā} again
\end{itemize}

\begin{tcolorbox}[breakable,colback=teal!4,colframe=teal!35!black,title={Before you move on}]
What does \hi{इकट्ठा} mean? (\textquotedblleft{}Gathered together\textquotedblright{}.) Say it once more.
\end{tcolorbox}
\section[\texorpdfstring{\hi{झाँकी} (jhā̃kī)}{झाँकी (jhā̃kī)}]{\hi{झाँकी} (jhā̃kī) — a tableau; a float in a procession}

\label{lesson:HI-C330-jhaki}



\begin{quote}
Before the new one: say the Hindi for merriment, then the Hindi for gathered together.
\end{quote}

\subsection*{You'll want to know: \hi{झाँकी}}
\textbf{\hi{झाँकी}} — \emph{jhā̃kī} — \textquotedblleft{}a tableau; a float in a procession\textquotedblright{}.

\begin{quote}
\hi{जुलूस} \hi{में} \hi{बच्चों} \hi{की} \hi{झाँकी} \hi{सबसे} \hi{सुंदर} \hi{थी।}
\end{quote}

\begin{quote}
\emph{julūs meṁ baccoṁ kī jhā̃kī sabse sundar thī.}
\end{quote}

\begin{quote}
The children's float was the most beautiful in the procession.
\end{quote}

Feminine.

\begin{cousinweb}
\textbf{\hi{झाँकी}} (\emph{jhā̃kī}) is built on \emph{jhā̃knā}, to peep, which CDIAL 5325 records as \emph{jhā̃khnā} under \emph{*jhaṅkh}: a glimpse, and then a scene set up for people to look at.
\end{cousinweb}

\begin{grammarlens}[title={What you've built}]
One more word for telling what happened.
\end{grammarlens}

\subsection*{Your turn}
\begin{itemize}
\raggedright
  \item \emph{Say it:} \emph{jhā̃kī}
  \item \emph{Say it:} \emph{jhā̃kī}, once more
  \item \emph{Say it:} \emph{julūs meṁ baccoṁ kī jhā̃kī sabse sundar thī.} — The children's float was the most beautiful in the procession.
  \item \emph{Recall:} say the Hindi for merriment, then the Hindi for gathered together, then say \emph{jhā̃kī} again
\end{itemize}

\begin{tcolorbox}[breakable,colback=teal!4,colframe=teal!35!black,title={Before you move on}]
What does \hi{झाँकी} mean? (\textquotedblleft{}A tableau; a float in a procession\textquotedblright{}.) Say it once more.
\end{tcolorbox}
\section[\texorpdfstring{\hi{झूलना} (jhūlnā)}{झूलना (jhūlnā)}]{\hi{झूलना} (jhūlnā) — to swing}

\label{lesson:HI-C330-jhulna}



\begin{quote}
Before the new one: say the Hindi for gathered together, then the Hindi for a tableau.
\end{quote}

\subsection*{You'll want to know: \hi{झूलना}}
\textbf{\hi{झूलना}} — \emph{jhūlnā} — \textquotedblleft{}to swing\textquotedblright{}.

\begin{quote}
\hi{बच्चे} \hi{झूले} \hi{पर} \hi{झूल} \hi{रहे} \hi{थे।}
\end{quote}

\begin{quote}
\emph{bacce jhūle par jhūl rahe the.}
\end{quote}

\begin{quote}
The children were swinging on the swing.
\end{quote}

Past (he …): \hi{वह} \hi{झूले} \hi{पर} \hi{झूला}, \emph{vah jhūle par jhūlā}.

\begin{cousinweb}
\textbf{\hi{झूलना}} (\emph{jhūlnā}) is inherited from \emph{*jhulyati}, a form CDIAL reconstructs (CDIAL 5406). \textbf{\hi{झूला}}, a swing, belongs with it: you \textbf{\hi{झूलते} \hi{हो}} on a \textbf{\hi{झूला}}.
\end{cousinweb}

\begin{grammarlens}[title={What you've built}]
One more word for telling what happened.
\end{grammarlens}

\subsection*{Your turn}
\begin{itemize}
\raggedright
  \item \emph{Say it:} \emph{jhūlnā}
  \item \emph{Say it:} \emph{jhūlnā}, once more
  \item \emph{Say it:} \emph{bacce jhūle par jhūl rahe the.} — The children were swinging on the swing.
  \item \emph{Recall:} say the Hindi for gathered together, then the Hindi for a tableau, then say \emph{jhūlnā} again
\end{itemize}

\begin{tcolorbox}[breakable,colback=teal!4,colframe=teal!35!black,title={Before you move on}]
What does \hi{झूलना} mean? (\textquotedblleft{}To swing\textquotedblright{}.) Say it once more.
\end{tcolorbox}
\section[\texorpdfstring{\hi{शोभायात्रा} (śobhāyātrā)}{शोभायात्रा (śobhāyātrā)}]{\hi{शोभायात्रा} (śobhāyātrā) — a festive procession}

\label{lesson:HI-C330-sobhayatra}



\begin{quote}
Before the new one: say the Hindi for a tableau, then the Hindi for to swing.
\end{quote}

\subsection*{You'll want to know: \hi{शोभायात्रा}}
\textbf{\hi{शोभायात्रा}} — \emph{śobhāyātrā} — \textquotedblleft{}a festive procession\textquotedblright{}.

\begin{quote}
\hi{शोभायात्रा} \hi{कितने} \hi{बजे} \hi{चलेगी}?
\end{quote}

\begin{quote}
\emph{śobhāyātrā kitne baje calegī?}
\end{quote}

\begin{quote}
What time will the procession set off?
\end{quote}

Feminine.

\begin{cousinweb}
\textbf{\hi{शोभायात्रा}} (\emph{śobhāyātrā}) joins Sanskrit \emph{śobhā}, \textquotedblleft{}splendour, beauty\textquotedblright{} (Monier-Williams), and \textbf{\hi{यात्रा}}, a journey or procession: a procession for show, with music and decorated floats. Like \textbf{\hi{जुलूस}}, it \textbf{\hi{निकलती} \hi{है}}.
\end{cousinweb}

\begin{grammarlens}[title={What you've built}]
One more word for telling what happened.
\end{grammarlens}

\subsection*{Your turn}
\begin{itemize}
\raggedright
  \item \emph{Say it:} \emph{śobhāyātrā}
  \item \emph{Say it:} \emph{śobhāyātrā}, once more
  \item \emph{Say it:} \emph{śobhāyātrā kitne baje calegī?} — What time will the procession set off?
  \item \emph{Recall:} say the Hindi for a tableau, then the Hindi for to swing, then say \emph{śobhāyātrā} again
\end{itemize}

\begin{tcolorbox}[breakable,colback=teal!4,colframe=teal!35!black,title={Before you move on}]
What does \hi{शोभायात्रा} mean? (\textquotedblleft{}A festive procession\textquotedblright{}.) Say it once more.
\end{tcolorbox}
\section[\texorpdfstring{\hi{भक्त} (bhakt)}{भक्त (bhakt)}]{\hi{भक्त} (bhakt) — a devotee}

\label{lesson:HI-C330-bhakt}



\begin{quote}
Before the new one: say the Hindi for to swing, then the Hindi for a festive procession.
\end{quote}

\subsection*{You'll want to know: \hi{भक्त}}
\textbf{\hi{भक्त}} — \emph{bhakt} — \textquotedblleft{}a devotee\textquotedblright{}.

\begin{quote}
\hi{हज़ारों} \hi{भक्त} \hi{मंदिर} \hi{के} \hi{बाहर} \hi{इकट्ठा} \hi{हुए।}
\end{quote}

\begin{quote}
\emph{hazāroṁ bhakt mandir ke bāhar ikaṭṭhā hue.}
\end{quote}

\begin{quote}
Thousands of devotees gathered outside the temple.
\end{quote}

Masculine or feminine, by the person meant.

\begin{cousinweb}
\textbf{\hi{भक्त}} (\emph{bhakt}) is Sanskrit \emph{bhakta}, \textquotedblleft{}devoted, attached\textquotedblright{} (Monier-Williams): a person devoted to a god.
\end{cousinweb}

\begin{grammarlens}[title={What you've built}]
One more word for telling what happened.
\end{grammarlens}

\subsection*{Your turn}
\begin{itemize}
\raggedright
  \item \emph{Say it:} \emph{bhakt}
  \item \emph{Say it:} \emph{bhakt}, once more
  \item \emph{Say it:} \emph{hazāroṁ bhakt mandir ke bāhar ikaṭṭhā hue.} — Thousands of devotees gathered outside the temple.
  \item \emph{Recall:} say the Hindi for to swing, then the Hindi for a festive procession, then say \emph{bhakt} again
\end{itemize}

\begin{tcolorbox}[breakable,colback=teal!4,colframe=teal!35!black,title={Before you move on}]
What does \hi{भक्त} mean? (\textquotedblleft{}A devotee\textquotedblright{}.) Say it once more.
\end{tcolorbox}
